% TOP_PROBLEM: {"id":30007627,"problem_number":"LOCAL-30007627","title":"Supply-chain resilience versus efficiency","category_id":21,"category":{"id":21,"name":"economics","display_name":"Economics","description":"Open economic theory statements, published research agendas, and proposed scoped specifications; question provenance and historical priority are qualified per record.","slug":"economics","order_index":21,"created_at":"2026-10-08T00:00:00Z"},"status":"research_question","external_url":null,"legacy_ids":[],"published":false,"statement":"In the explicitly specified setting: What diversification, inventory and sourcing policies optimally trade lower expected costs against rare disruption risks? The mathematical statement above fixes the target, assumptions and scope of this catalogue formulation.","economics":{"original_id":116,"field":"Trade, globalization and geoeconomics","kind":"design","formalization_basis":"adapted_research_question","source_status":"research_agenda","priority_status":"earliest_found","priority_note":"This is the earliest explicit statement verified in this audit; first priority for the full catalogue formulation is not established.","earliest_verified_reference":{"authors":["Bank of England"],"title":"Bank of England Agenda for Research, 2025–2028","year":2026,"url":"https://www.bankofengland.co.uk/research/bank-of-england-agenda-for-research","doi":null,"locator":"Interconnected World / Reconfigured trade, paragraph2; Priority topics for2026, item2","evidence_summary":"Requests complex supply-chain resilience and vulnerability research."},"current_reference":{"authors":["Bank of England"],"title":"Bank of England Agenda for Research, 2025–2028","year":2026,"url":"https://www.bankofengland.co.uk/research/bank-of-england-agenda-for-research","doi":null,"locator":"Interconnected World / Reconfigured trade, paragraph2; Priority topics for2026, item2","evidence_summary":"Requests complex supply-chain resilience and vulnerability research."},"formalization_note":"The source asks resilience and supply-chain vulnerability questions. A robust planner objective and uncertainty set are proposed, not a canonical universal open theorem."},"statusQualification":"Published question or research agenda verified at the cited locator. The mathematical specification is adapted for this catalogue and includes additional assumptions; source publication does not establish first-ever priority or certify this exact model remains unsolved. The source asks resilience and supply-chain vulnerability questions. A robust planner objective and uncertainty set are proposed, not a canonical universal open theorem.","statusReviewedAt":"2026-10-08"}
% FIRST_POSTED: 2026-10-08
% ENTRY_KIND: statement_only
% !TeX program = lualatex
\documentclass[11pt]{article}
\usepackage{fontspec}
\usepackage[a4paper,margin=25mm]{geometry}
\usepackage{amsmath,amssymb,mathtools}
\usepackage{xurl}
\usepackage[hidelinks]{hyperref}
\newcommand{\sref}[2]{\hyperlink{source:#1:#2}{[S#2]}}
\setlength{\emergencystretch}{3em}
\begin{document}
% problemId:30007627
\section*{Supply-chain resilience versus efficiency}\label{30007627}
\subsection{Definitions and mathematical statement}
Consider firms choosing supplier shares \(w\), inventories \(b\) and substitution capabilities \(s\) before uncertain input disruptions. Feasible configurations belong to \(A\), defined by production technology, working-capital and contracting constraints. A disruption state \(\omega\) specifies supplier availability and input costs; delivered output and consumption follow from equilibrium production and price responses. Let \(C(a)\) be normal resource costs and \(L(a,\omega)\) disruption losses for configuration \(a=(w,b,s)\). For a publicly specified set \(\mathcal P\) of shock distributions, seek
\[a^*\in\arg\min_{a\in A}\{C(a)+\sup_{P\in\mathcal P}E_P[L(a,\omega)]\}.\]
The task is to identify the efficiency-resilience frontier and whether private firm choices attain the social optimum when shortages spill through production networks. Estimate switching costs, substitutability and correlated supplier risks rather than assuming suppliers fail independently. Compare diversification, inventory and domestic sourcing using the same loss convention. Include endogenous supplier prices and entry when diversification is adopted widely. Report how optimal decisions vary with upstream position and shock severity, and give bounds where rare-event probabilities are poorly measured. Domestic sourcing is not defined as resilient by assumption; concentrated domestic supply can face common shocks.

\par\textbf{Formulation and scope.} The source asks resilience and supply-chain vulnerability questions. A robust planner objective and uncertainty set are proposed, not a canonical universal open theorem.

\par\textbf{Open-question provenance.} An explicit question or research agenda was verified at the source locator. This does not by itself certify that every later version remains unresolved \sref{30007627}{1}.

\par\textbf{Historical priority.} This is the earliest explicit statement verified in this audit; first priority for the full catalogue formulation is not established.

\subsection{Short English statement}
In the explicitly specified setting: What diversification, inventory and sourcing policies optimally trade lower expected costs against rare disruption risks? The mathematical statement above fixes the target, assumptions and scope of this catalogue formulation.

\subsection{Sources}
\begin{itemize}\raggedright
\item[\textnormal{[S1]}]\hypertarget{source:30007627:1}{} Bank of England. 2026. Bank of England Agenda for Research, 2025–2028. Interconnected World / Reconfigured trade, paragraph2; Priority topics for2026, item2.
\url{https://www.bankofengland.co.uk/research/bank-of-england-agenda-for-research}.
{\small\textit{Earliest verified question/agenda reference; historical priority qualified below. Requests complex supply-chain resilience and vulnerability research.}}
\end{itemize}
\end{document}
